\section{真题代练}

\subsection{现代文阅读 I：信息类}

\subsubsection{真题 1（2024 新课标Ⅰ卷 · 信息类阅读）}

阅读下面的文字，完成 1--4 题。

（材料略，此处为示例结构）

\textbf{解题思路}：

\begin{enumerate}
  \item 第 1 题（信息筛选）：定位原文第 2 段，选项 C 将「可能」偷换为「一定」
  \item 第 2 题（推断判断）：选项 A 曲解文意，原文未涉及因果
  \item 第 3 题（论证分析）：选项 D 论证方法判断错误，此处为举例论证而非对比论证
  \item 第 4 题（概括简答）：总分总结构，先用一句话概括观点，再分条列举理由
\end{enumerate}

\subsubsection{真题 2（2023 新课标Ⅰ卷）}

\textbf{考查重点}：筛选整合 + 图表信息转换

\textbf{失分点}：图表与文字信息不对应、范围词理解偏差

\subsubsection{真题 3（2022 新课标Ⅰ卷）}

\textbf{考查重点}：分析论证特点 + 评价作者观点

\textbf{答题提示}：论证特点从方法、结构、语言三方面展开

\subsection{现代文阅读 II：文学类}

\subsubsection{真题 1（2024 新课标Ⅰ卷 · 文学类阅读）}

\textbf{文体}：报告文学

\textbf{高频考点}：
\begin{itemize}
  \item 第 6 题：理解句子含义（联系上下文 + 表层深层）
  \item 第 7 题：艺术特色判断（注意人称变化的作用）
  \item 第 8 题：分析叙述视角（第一人称增强真实感）
  \item 第 9 题：探究标题意蕴（表层含义 + 深层象征）
\end{itemize}

\subsubsection{真题 2（2023 新课标Ⅰ卷）}

\textbf{文体}：散文

\textbf{考查重点}：句段作用分析 + 语言风格赏析

\subsubsection{真题 3（2021 新课标Ⅰ卷）}

\textbf{文体}：小说

\textbf{考查重点}：人物形象概括 + 情节作用 + 主题探究

\subsection{文言文阅读}

\subsubsection{真题 1（2024 新课标Ⅰ卷）}

\textbf{选文}：《通鉴纪事本末》

\textbf{断句}（10 题）：注意专有名词与句末语气词的识别

\textbf{翻译}（13 题）：
\begin{itemize}
  \item 得分点 1：「……」，注意被动句式
  \item 得分点 2：「……」，注意词类活用
  \item 整体句式：状语后置句，需调整语序
\end{itemize}

\subsubsection{真题 2（2023 新课标Ⅰ卷）}

\textbf{选文}：《韩非子》/《说苑》等诸子散文

\textbf{考查重点}：论证说理特点 + 对比手法在论说中的作用

\subsubsection{真题 3（2022 新课标Ⅰ卷）}

\textbf{选文}：《战国策》体史书

\textbf{考查重点}：人物劝说策略 + 事实与观点的区分

\subsection{古代诗歌鉴赏}

\subsubsection{真题 1（2024 新课标Ⅰ卷）}

\textbf{诗歌类型}：唐诗（边塞诗/送别诗）

\textbf{选择题设误}：意象解读错误（将「羌笛」的哀怨说成欢快）

\textbf{简答题}：赏析某一联的表达效果

\[
\text{答案结构}：手法（视听结合/借景抒情）+ 内容描写 + 情感表达
\]

\subsubsection{真题 2（2023 新课标Ⅰ卷）}

\textbf{诗歌类型}：宋诗（哲理诗）

\textbf{考查重点}：比较阅读 + 诗中哲理的理解

\subsubsection{真题 3（2021 新课标Ⅰ卷）}

\textbf{诗歌类型}：唐诗（咏物诗）

\textbf{考查重点}：托物言志手法 + 物象与品格的联系

\subsection{语言文字运用}

\subsubsection{真题 1（2024 新课标Ⅰ卷）}

\textbf{考查形式}：一材两题

\begin{itemize}
  \item 第 1 题：成语填空（近义辨析 + 语境判断）
  \item 第 2 题：病句修改（语序不当 + 成分残缺）
\end{itemize}

\subsubsection{真题 2（2023 新课标Ⅰ卷）}

\textbf{考查形式}：语句复位 + 修辞赏析

\textbf{语句复位}：注意上下文话题一致性与句式对称

\textbf{修辞赏析}：比喻（本体+喻体+相似点）或拟人（赋予人的特征+生动形象）

\subsubsection{真题 3（2022 新课标Ⅰ卷）}

\textbf{考查形式}：标点符号作用 + 句式变换

\textbf{标点}：破折号（解释说明/话题转换/强调）与引号（特殊含义/引用/讽刺）

\subsection{写作真题拆解与示范}

\subsubsection{真题 1（2024 新课标Ⅰ卷）——科技发展与人文关怀}

\textbf{材料核心}：人工智能时代，人们越来越多地依赖技术解决问题。「随着互联网的普及、人工智能的应用，越来越多的问题能很快得到答案。那么，我们的问题是否会越来越少？」

\textbf{审题拆解}：

\begin{enumerate}
  \item \textbf{圈关键词}：互联网 + 人工智能 + 答案变多 + 问题变少？
  \item \textbf{判断关系}：这是一个「表面现象 vs 深层本质」的思辨话题
  \item \textbf{确立立场}：好的回答是——问题不会减少，反而会更多、更深
  \item \textbf{提炼中心论点}：\textbf{技术给我们答案，但真正的思考才刚刚开始}
\end{enumerate}

\textbf{立意方向 + 万金油转化}：

\begin{center}
\begin{tabular}{c|c}
\textbf{方向} & \textbf{如何引向爱国/青年奉献} \\ \hline
思辨方向：答案越多，问题越深 & 分论点 3：当代青年在技术浪潮中，更需以家国情怀为坐标系，思考「技术向善」这一时代命题——这本身就是一种奉献 \\
务实方向：技术解放人力，专注创造 & 分论点 3：被技术解放的青年，应把精力投入到国家最需要的地方——乡村振兴、科技攻关、文化传承 \\
\end{tabular}
\end{center}

\textbf{全文框架示范}：

\begin{itemize}
  \item[\textbf{开头}] 引用现象 + 设问：「ChatGPT 秒回所有问题，是不是意味着人类不再需要提问了？」$\to$ 中心论点
  \item[\textbf{分论点 1}] 知识性问题的减少，倒逼我们提出更有价值的问题（个人层面）
  \item[\textbf{分论点 2}] AI 能回答「是什么」，却不能回答「为什么」和「应该怎样」（社会层面）
  \item[\textbf{分论点 3}] \textbf{青年面对 AI 时代，当以独立思考为基，以家国情怀为舵}（国家层面——青年奉献）
  \item[\textbf{结尾}] 升华：问题不是越来越少，而是越来越高——这是文明的进步，更是青年的使命
\end{itemize}

\textbf{可用素材}：屠呦呦（不依赖 AI，从古籍中发现青蒿素）、中国航天（自主创新突破卡脖子技术）、ChatGPT 的伦理困境

\subsubsection{真题 2（2023 新课标Ⅰ卷）——个人成长与时代担当}

\textbf{材料核心}：好的故事，可以改变一个人的命运，可以展现一个民族的形象……故事是有力量的。

\textbf{审题拆解}：

\begin{enumerate}
  \item \textbf{圈关键词}：好故事 + 改变命运 + 展现民族形象 + 力量
  \item \textbf{判断关系}：递进关系——从个人到民族
  \item \textbf{确立立场}：既要听好故事，更要做写故事的人
  \item \textbf{中心论点}：\textbf{好故事的力量，不仅在于打动人心，更在于激励我们把个人故事写进时代的叙事中}
\end{enumerate}

\textbf{万金油转化}：

「听好故事」$\to$ 感受榜样的力量（爱国故事）$\to$「做好故事的续写者」$\to$ 拿过接力棒，用青春书写新时代的中国故事

\textbf{全文框架示范}：

\begin{itemize}
  \item[\textbf{开头}] 排比引入：「从『为中华之崛起而读书』到『清澈的爱，只为中国』，一代代中国青年的故事，构成了我们这个民族最动人的精神底色。」$\to$ 中心论点
  \item[\textbf{分论点 1}] 好故事塑造个人品格（举例：读《红岩》懂得坚韧）
  \item[\textbf{分论点 2}] 好故事凝聚民族精神（举例：抗疫故事、航天故事）
  \item[\textbf{分论点 3}] \textbf{当代青年不能只做故事的听众，更要做故事的创造者}
  \item[\textbf{结尾}] 呼吁式结尾：「我们是什么样子，中国故事就是什么样子。」
\end{itemize}

\subsubsection{真题 3（2022 新课标Ⅰ卷）——传承与创新}

\textbf{材料核心}：围棋中有「本手、妙手、俗手」三个术语——本手是基础，妙手是创造，俗手是忽视基础的貌似巧妙。

\textbf{审题拆解}：

\begin{enumerate}
  \item \textbf{圈关键词}：本手（基础）+ 妙手（创新）+ 俗手（好高骛远）
  \item \textbf{判断关系}：对立统一——基础与创新的辩证
  \item \textbf{确立立场}：没有扎实的本手，就没有真正的妙手
  \item \textbf{中心论点}：\textbf{任何领域的卓越，都建立在日复一日的「笨功夫」之上}
\end{enumerate}

\textbf{万金油转化}：

本手 = 脚踏实地的积累（黄文秀走遍每一户人家）
妙手 = 关键时刻的创新突破（航天青年的技术攻坚）
把「下好本手、追求妙手」升华为「青年扎根基层、奉献祖国」

\textbf{全文框架示范}：

\begin{itemize}
  \item[\textbf{开头}] 围棋引入 + 类比到人生：「人生如棋，有人追求妙手，有人不屑本手。然而真正的智慧在于——没有本手的千锤百炼，妙手不过是空中楼阁。」
  \item[\textbf{分论点 1}] 本手是根基：任何成就都离不开扎实积累（袁隆平数十年田间观察）
  \item[\textbf{分论点 2}] 妙手是升华：创新建立在对基础规律深刻理解之上（屠呦呦从古籍中找到灵感）
  \item[\textbf{分论点 3}] \textbf{当代青年既要有脚踏实地的定力（本手），也要有敢为人先的锐气（妙手），归根结底是要把个人奋斗融入国家发展的棋局之中}
  \item[\textbf{结尾}] 总结升华：「人生这盘棋，下好本手是态度，追求妙手是境界。而把每一步都走在家国需要的大棋盘上，才是我们这一代青年最精彩的落子。」
\end{itemize}

\subsubsection{真题 4（2021 新课标Ⅰ卷）——追求理想与自身修养}

\textbf{材料核心}：汉代扬雄以射箭为喻——「修身以为弓，矫思以为矢，立义以为的，奠而后发，发必中矣。」（修身是弓，端正思想是箭，确立道义是靶心）

\textbf{审题拆解}：

\begin{enumerate}
  \item \textbf{圈关键词}：修身 + 矫思 + 立义 = 实现理想
  \item \textbf{判断关系}：条件关系——修身、矫思、立义是实现理想的三个前提
  \item \textbf{确立立场}：理想不是空想，而是修身立德后的水到渠成
  \item \textbf{中心论点}：\textbf{真正的理想追求，始于修身、成于正思、终于立义}
\end{enumerate}

\textbf{万金油转化}：

立义 = 把个人理想与国家需要相结合
结尾落脚：「当代青年当以国家之『义』为靶心，把个人理想的箭矢射向时代最需要的地方。」

\textbf{可用素材}：钟南山（修身——医术精湛）、张桂梅（矫思——坚定的教育信念）、边防战士（立义——以国为靶心）

\subsubsection{真题 5（2020 新课标Ⅰ卷）——历史人物评说}

\textbf{材料核心}：齐桓公、管仲、鲍叔三人的故事，班级读书会请你就其中感触最深的人物写一篇发言稿。

\textbf{审题拆解}：

\begin{enumerate}
  \item \textbf{关键词}：齐桓公（宽容大度）+ 管仲（才能报国）+ 鲍叔（知人荐贤）
  \item \textbf{任务指令}：\textbf{发言稿}（注意格式）+ 选感触最深的一个
  \item \textbf{中心论点}（选鲍叔为例）：\textbf{鲍叔最令我感动的，不是他的能力，而是他的格局——一种把国家利益置于个人得失之上的大格局}
\end{enumerate}

\textbf{万金油转化}：

鲍叔的格局 = 当代青年的奉献精神
结尾升华：「鲍叔的格局穿越千年依然动人，因为它昭示着一个朴素的真理——一个人的价值不在于他占据什么位置，而在于他为更广阔的事业贡献了什么。这正是我们今天所呼唤的奉献精神。」

\subsection{写作素材积累引导}

\subsubsection{素材积累的「三三制」}

\textbf{三个维度}：同一人物可从不同角度使用

\begin{center}
\begin{tabular}{c|c|c}
\textbf{人物} & \textbf{可以用在什么主题} & \textbf{一句话素材} \\ \hline
袁隆平 & 爱国/奉献/坚持/科学精神/脚踏实地 & 「一粒种子改变世界，他用一生践行了一个誓言：把饭碗掌握在中国人自己手里。」 \\
黄文秀 & 青春奉献/理想/担当/平凡伟大 & 「从北师大硕士到百坭村第一书记，她走过的两万五千里，是新时代的长征路。」 \\
屠呦呦 & 创新/坚持/文化自信/科技报国 & 「从古籍中捕捉灵感，用 190 次失败换来一株青蒿拯救全球数百万人。」 \\
南仁东 & 爱国/坚守/科学精神/追求卓越 & 「24 年只做一件事——把中国天眼建在贵州的大山里，让中国看到更远的宇宙。」 \\
樊锦诗 & 文化传承/坚守/热爱/奉献 & 「从北大才女到大漠守护者，她用一生守住了敦煌，也守住了我们的文化根脉。」 \\
张桂梅 & 奉献/教育/信念/女性力量 & 「她把 1800 多个女孩送出大山，自己却一身病痛。她说：只要还有一口气，就要站在讲台上。」 \\
\end{tabular}
\end{center}

\textbf{三个层次}：一个人物要记住三个层次
\begin{enumerate}
  \item \textbf{一句话概括}（用在开头/结尾快速例证）
  \item \textbf{一个细节}（用在主体段展示深度）
  \item \textbf{一句名言/数据}（增强说服力）
\end{enumerate}

\textbf{示范}（以袁隆平为例）：
\begin{enumerate}
  \item 一句话：「袁隆平用一粒种子改变了世界。」
  \item 细节：他 90 岁还在下田，每天第一件事就是去试验田走一圈。他说：「我不在试验田，就在去试验田的路上。」
  \item 数据：杂交水稻累计种植面积超过 6 亿公顷，增产粮食近 8 亿吨。
\end{enumerate}

\subsubsection{如何用一个素材套用多个主题}

\textbf{素材万用公式}：

\[
\text{人物事迹（客观事实）} + \text{主题关键词连接} = \text{贴合任意主题}
\]

\textbf{示范——以「黄文秀」为例}：

\begin{center}
\begin{tabular}{c|c}
\textbf{作文主题} & \textbf{黄文秀素材的切入角度} \\ \hline
青春与理想 & 她放弃了城市的繁华，选择了最艰苦的基层——这才是青春该有的样子 \\
奉献与担当 & 她驻村满一年，汽车里程正好两万五千里——用生命践行了新时代的长征精神 \\
平凡与伟大 & 她做的都是走访、修路、找销路这些「小事」，但积小流成江海 \\
选择与价值 & 北师大的学历不是她向上爬的阶梯，而是她向下扎根的底气 \\
坚持与信念 & 暴雨之夜她驱车回村，只因担心村民的防汛——直到生命最后一刻 \\
\end{tabular}
\end{center}

\textbf{关键技巧}：同一个素材，开头用一句话引出 + 中间用细节展开 + 结尾用金句收束，同一个素材可以出现在分论点 1/2/3 的不同位置。

\subsubsection{每周积累计划}

\begin{center}
\begin{tabular}{c|c|c}
\textbf{星期} & \textbf{任务} & \textbf{时长} \\ \hline
周一 & 精读 1 个人物素材（记 1 个细节 + 1 句名言） & 15 分钟 \\
周三 & 精读 1 个时事素材（关注人民日报 / 新华社新媒体） & 15 分钟 \\
周五 & 看 1 篇高考优秀范文（拆解结构 + 记 3 个好句子） & 20 分钟 \\
周末 & 练习 1 段 200 字写作（用本周素材写一个主体段） & 30 分钟 \\
\end{tabular}
\end{center}

\textbf{推荐获取渠道}：
\begin{itemize}
  \item 《人民日报》评论版、「人民日报评论」微信公众号
  \item 「央视新闻」人物专栏、《感动中国》获奖人物
  \item 高考满分作文集（重点看开头结尾和结构安排）
\end{itemize}
